% Principal GDZ PPN236018647,native476–479/printed469–472;HKUST477–480.
\clearpage\SourceLeaf{476}\NoSourceHead
\pdfbookmark[1]{Mannigfaltigkeiten von n Dimensionen}{manifolds-n}
\HistoricTitle{[MANNIGFALTIGKEITEN VON \(n\) DIMENSIONEN.]}
\begin{center}[I.]\end{center}
\HistoricTitle{ÜBER DIE ERMITTELUNG DES KLEINSTEN WERTHES\\ FÜR DAS MAASS DES ZWANGES.}
\HistoricCitation{[\textsc{A. Ritter}, \textit{Ueber das Princip des kleinsten Zwanges}, Inaugural-Dissertation, Göttingen 1853, S. 20—23.]}
\indent \{Die Regel, welche die Mathematik für das Aufsuchen der grössten und kleinsten Werthe einer »stetigen« Funktion angiebt, ist bekanntlich die: man variire die Funktion und setze die Coefficienten der von einander unabhängigen Variationen gleich Null; die Auflösung der hierdurch erhaltenen Gleichungen ergiebt die gesuchten Werthe der unabhängigen Veränderlichen. Die Gleichung \(dZ=0\) würde daher unsere Aufgabe sofort zum Abschluss bringen unter der Voraussetzung, dass \(Z\) eine stetige Funktion ist; sie ist der Ausdruck für das \textsc{Lagrange}sche Princip, also identisch mit derjenigen, welche aus der Verbindung des Princips der virtuellen Geschwindigkeiten mit dem \textsc{Dalembert}schen Princip sich ergiebt. Sie enthält in der That die allgemeinen Gesetze der Bewegung eines Systems materieller Punkte für alle die Fälle, wo die Bedingungen, denen die Bewegung unterworfen ist, sich durch »Gleichungen« ausdrücken lassen.

\indent Anders verhält es sich indessen, wenn unter den gegebenen äusseren Beschränkungen auch solche vorkommen, welche nur durch »Ungleichungen«, z. B. von der Form
\[
f(\xi,\eta,\zeta,.\ .\ .)\geqq 0
\]
\noindent auszudrücken sind. In diesem Falle wird die Funktion \(Z\) unstetig, und der Anwendung jener Regel muss dann eine Untersuchung desjenigen Theils ihrer
\clearpage\SourceLeaf{477}\ArithmeticHead{470}{L}{GEOMETRIE. NACHLASS.}
\noindent Werthen-Reihe vorangehen, welcher bei der vorliegenden Frage in Betracht kommt.

\indent Mit Hülfe jeder gegebenen Bedingungs-Gleichung lässt sich Eine der veränderlichen Grössen in der Funktion \(Z\) eliminiren. Nennen wir den Inbegriff aller Werthe einer von \(N\) veränderlichen Grössen abhängigen Funktion für alle Combinationen dieser Grössen von \(-\infty\) bis \(+\infty\) einen \textls[120]{analytischen Raum von} \(N\) \textls[120]{Dimensionen}\footnote{Die hierdurch wesentlich erleichterte Ausdrucksweise möge diese Einkleidung entschuldigen.}, so wird die Berücksichtigung jeder »Gleichung«, die Wirkung haben, dass sie die Anzahl der Dimensionen des analytischen Raumes der Funktion \(Z\) um Eine vermindert. Jede »Ungleichung« dagegen beschränkt das Gebiet der Veränderlichkeit, ohne eine Dimension hinwegzunehmen, sie theilt den ganzen analytischen Raum der Funktion in zwei Theile, deren Einer als »möglich«, während der andere als »unmöglich« betrachtet werden soll.

\indent Unsere Aufgabe schreibt nun vor: es soll unter allen »möglichen« Werthen-Combi\-nationen diejenige gesucht werden, welche die Funktion \(Z\) zu einem Minimum macht. Unsere nächste Sorge wird daher die sein: die Ausdehnung des »möglichen« Gebietes in dem analytischen Raum jener Funktion kennen zu lernen. Es sei nun durch Berücksichtigung sämmtlicher »Gleichungen« die Anzahl seiner Dimensionen bis auf \(n\) vermindert, und die Anzahl der noch zu berücksichtigenden »Ungleichungen« sei \(=\mu\); dann wird, wenn wir den ganzen analytischen Raum von \(n\) Dimensionen mit \(\mathfrak{R}_n\) bezeichnen, die Wirkung der Ungleichungen die sein: dass sie das mögliche Gebiet auf ein begrenztes Stück des Raumes \(\mathfrak{R}_n\) einschränkt, und zwar wird dieses Stück, welches wir \(\mathfrak{r}_n\) nennen, immer noch einen Raum von \(n\) Dimensionen bilden.

\indent Es sind nun zwei Fälle möglich: entweder liegt der Punkt, welcher die gesuchte Werthen-Combination repräsentirt, »innerhalb« des Raumes \(\mathfrak{r}_n\), oder er liegt in der Begrenzung desselben. Um uns hierüber Aufschluss zu verschaffen, suchen wir zunächst denjenigen Punkt des Raumes \(\mathfrak{R}_n\) auf, welcher das absolute Minimum der Funktion \(Z\) repräsentirt, d. h. wir suchen nach der bekannten Regel ihr Minimum, ohne die Ungleichungen zu berücksichtigen. Findet sich dann, dass die Substitution der gefundenen Werthe alle Unglei-
\clearpage\SourceLeaf{478}\ArithmeticHead{471}{R}{MANNIGFALTIGKEITEN VON \(n\) DIMENSIONEN.}
\noindent chungen positiv macht, so ist die Aufgabe gelöst, und es ist so gut als wären gar keine Ungleichungen vorhanden gewesen.

\indent Im entgegengesetzten Falle ist der Punkt in der Begrenzung des Stückes \(\mathfrak{r}_n\) zu suchen, und die anzuwendende Methode wird im Allgemeinen folgende sein. Vorausgesetzt: dass die Ungleichungen von einander unabhängig sind, d. h. dass keine von ihnen in einer oder mehreren der anderen enthalten ist, werden die Grenzen des Stücks \(\mathfrak{r}_n\) durch \(\mu\) Räume von \(n-1\) Dimensionen gebildet, und je \(n\) von diesen Grenz-Räumen schneiden sich in Einem Punkte; oder — wenn \(n>\mu\) — alle Grenz-Räume haben Einen Raum von \(n-\mu\) Dimensionen gemeinschaftlich. Der erste Schritt ist nun der: dass wir Einen der Durchschnittspunkte (im ersten Falle) oder Einen Punkt des Durchschnitts-Raums (im letzteren) aufsuchen; dieser Punkt bildet den Ausgangspunkt für unsere Untersuchung. Wir finden ihn, indem wir irgend ein Werthen-System aufsuchen, welches \(n\) von den Ungleichungen \(=\) Null und die übrigen positiv, oder — wenn \(n>\mu\) — welches alle \(\mu\) Ungleichungen \(=\) Null werden lässt. Von diesem Punkte ausgehend untersuchen wir weiter: in welcher Richtung wir von ihm fortschreiten müssen, um den Werth der Funktion \(Z\) am schnellsten abnehmen zu lassen, oder auf welche Weise eine gegebene Aenderung von \(Z\) durch möglich kleinste Aenderung des Werthen-Systems zu erreichen ist, wobei als Maass der letzteren die Summe der Quadrate der einzelnen Aenderungen betrachtet werden kann. In der gefundenen Richtung lassen wir dann den Punkt sich ändern, bis entweder derjenige Punkt, in welchem die Funktion ein Minimum wird, erreicht ist, oder bis eine der vorher positiven Ungleichungen aufhört erfüllt zu sein, d. h. Null wird. Im ersteren Falle ist die Aufgabe gelöst, im letzteren bildet der gefundene Punkt den Ausgangspunkt für eine ähnliche Untersuchung, welche ihrerseits dann wieder zu einem anderen Punkte führt, der entweder selbst das Minimum repräsentirt, oder wieder als Ausgangspunkt für einen dritten Schritt dient — kurz die ganze Operation wird so lange fortgesetzt, bis man zu einem Werthen-Systeme gelangt, welches die Eigenschaft besitzt, dass jede »mögliche« Aenderung desselben mit einem Grösserwerden der Funktion \(Z\) verbunden ist. Die Grenz-Räume, in deren gemeinschaftlichem Durchschnittsraum der endlich gefundene Punkt liegt, repräsentiren die »wirksamen« Ungleichungen, d. h. diejenigen, welche als Gleichungen erfüllt werden.
\clearpage\SourceLeaf{479}\ArithmeticHead{472}{L}{GEOMETRIE. NACHLASS.}
\indent Die Aufgabe würde daher vollkommen die nämliche gewesen sein, wenn diese wirkenden Ungleichungen von vornherein als Gleichungen, die übrigen Ungleichungen aber gar nicht gegeben wären. Die mögliche Veränderlichkeit der Funktion \(Z\) wäre dann schon durch die Natur der Aufgabe auf die richtige Anzahl der »unabhängigen« Veränderlichen beschränkt und innerhalb dieser Beschränkung eine stetige gewesen; und in diesem Falle würde also der Anwendung des \textsc{Lagrange}schen Princips oder der hiernach modificirten Gleichung \(dZ=0\) Nichts entgegengestanden haben; während sonst, in der Ungewissheit über die wirkenden Ungleichheiten, oder über die wirklich unabhängigen Veränderlichen, jenes Princip, um allgemein gültig zu sein, nur in der Form \(dZ\geqq 0\) ausgesprochen werden dürfte, also in Form einer Ungleichung, eines indirecten Urtheils. Dem Sinne nach würde es in dieser Fassung mit dem \textsc{Gauss}schen Princip übereinstimmen, ohne jedoch die Eleganz der Form und die unmittelbare Evidenz des letzteren zu erreichen.\}
\ClosingRule
